\par\Needspace{8\baselineskip}\bigskip
\pdfbookmark[3]{Exposition, 8}{pallas-II-article8}
\begin{center}8.\end{center}
\noindent\hspace*{\parindent}Nous nous servirons dans cette recherche des mêmes signes, que nous avons introduits dans la première section, non seulement pour les élémens proprement dits mais aussi pour les autres quantités auxiliaires. Ainsi p. e. \(A\) exprimera la valeur de \(\dfrac{y\mathrm{d}z-z\mathrm{d}y}{\mathrm{d}t}\) désormais variable; et on conçoit que toutes les équations finies entre ces quantités \(A,B,C,F,G,H,v,E\) etc., les élémens et les coordonnées \(x,y,z,r\) avec leurs différentielles \(\dfrac{\mathrm{d}x}{\mathrm{d}t}\mathpunct{\raisebox{.45ex}{,}}\dfrac{\mathrm{d}y}{\mathrm{d}t}\mathpunct{\raisebox{.45ex}{,}}\dfrac{\mathrm{d}z}{\mathrm{d}t}\mathpunct{\raisebox{.45ex}{,}}\dfrac{\mathrm{d}r}{\mathrm{d}t}\) seront indépendantes de cette variabilité. Il sera donc permis, de différentier toute équation finie entre les coordonnées \(x,y,z,r\), les élémens et les autres quantités, qui sont constantes dans l’hypothèse elliptique et qu’on peut comprendre aussi sous le nom d’élémens, en n’ayant égard qu’aux différentielles des coordonnées, et puisqu’en retranchant ce résultat de la différentielle complète il reste la différentielle de l’équation proposée, en ne traitant comme variables que les élémens seuls, il est évident, que cette nouvelle différentiation est aussi légitime. Ce théorème remarquable peut quelques fois servir à abréger les calculs, cependant on voit qu’on peut aussi toujours s’en passer.
